% !TEX encoding = UTF-8 Unicode
\documentclass[12pt,a4paper]{article}

% ── Language & encoding ──────────────────────────────────────────────────────
\usepackage[utf8]{inputenc}
\usepackage[T1]{fontenc}
\usepackage[spanish]{babel}

% ── Typography ────────────────────────────────────────────────────────────────
\usepackage{lmodern}
\usepackage{microtype}

% ── Page layout ──────────────────────────────────────────────────────────────
\usepackage[top=2.5cm, bottom=2.5cm, left=3cm, right=2.5cm]{geometry}
\usepackage{setspace}
\onehalfspacing

% ── Tables ───────────────────────────────────────────────────────────────────
\usepackage{booktabs}
\usepackage{tabularx}
\usepackage{array}

% ── Figures & floats ─────────────────────────────────────────────────────────
\usepackage{graphicx}
\usepackage{float}

% ── Lists ────────────────────────────────────────────────────────────────────
\usepackage{enumitem}

% ── Colors ───────────────────────────────────────────────────────────────────
\usepackage[dvipsnames]{xcolor}

% ── Bibliography ─────────────────────────────────────────────────────────────
% Load natbib BEFORE hyperref to avoid citation-link conflicts
\usepackage[numbers]{natbib}

% ── Hyperlinks ───────────────────────────────────────────────────────────────
\usepackage[colorlinks=true,
            linkcolor=blue!60!black,
            citecolor=blue!60!black,
            urlcolor=blue!70!black]{hyperref}

% ── Header / footer ──────────────────────────────────────────────────────────
\usepackage{fancyhdr}
\setlength{\headheight}{14pt}
\pagestyle{fancy}
\fancyhf{}
\rhead{\small Clasificaci\'{o}n de Afirmaciones Pol\'{i}ticas --- UPM ABID}
\cfoot{\thepage}
\renewcommand{\headrulewidth}{0.4pt}

% ── Caption style ────────────────────────────────────────────────────────────
\usepackage{caption}
\captionsetup{font=small, labelfont=bf}

% ─────────────────────────────────────────────────────────────────────────────
\begin{document}

% ── Title page ───────────────────────────────────────────────────────────────
\begin{titlepage}
  \centering
  \vspace*{2cm}
  {\Large\bfseries Universidad Politécnica de Madrid\\[0.4em]
  Escuela Técnica Superior de Ingeniería Informática\par}
  \vspace{1.5cm}
  \rule{\linewidth}{0.5pt}\\[0.5em]
  {\LARGE\bfseries Clasificación Automática de\\[0.3em]
  Afirmaciones Políticas\\[0.3em]
  mediante Procesamiento del Lenguaje Natural\par}
  \vspace{0.5em}\rule{\linewidth}{0.5pt}
  \vspace{1.2cm}
  {\large Competencia Kaggle --- Asignatura ABID\par}
  \vspace{1.5cm}
  {\normalsize
  \renewcommand{\arraystretch}{1.4}
  \begin{tabular}{rl}
    \textbf{Tarea:}       & Clasificación binaria de afirmaciones políticas \\
    \textbf{Dataset:}     & 8.950 muestras de entrenamiento + 3.860 de test \\
    \textbf{Métrica:}     & Macro F1 Score \\
    \textbf{Plataforma:}  & Kaggle \\
  \end{tabular}\par}
  \vspace{1.5cm}
  {\large\bfseries Grupo 1\par}
  \vspace{0.6cm}
  {\normalsize
  \begin{tabular}{c}
    Marla Gonz\'{a}lez Escobar               \\[0.3em]
    Bel\'{e}n Mar\'{i}a Iniesta Tejera       \\[0.3em]
    Silvia Rodr\'{i}guez Hern\'{a}ndez       \\[0.3em]
    Rosali Rodriguez Ure\~{n}a               \\[0.3em]
    Marilyn Salazar Mart\'{i}nez             \\[0.3em]
    Iv\'{a}n Omar Angeles Mart\'{i}nez       \\
  \end{tabular}\par}
  \vfill
  {\normalsize Mayo de 2025\par}
\end{titlepage}

\tableofcontents
\newpage

% ─────────────────────────────────────────────────────────────────────────────
\section{Introducción}
% ─────────────────────────────────────────────────────────────────────────────

La detección automática de afirmaciones falsas es una tarea cada vez más relevante dentro del Procesamiento del Lenguaje Natural, especialmente en contextos políticos, donde la desinformación puede influir de forma significativa en la opinión pública. Este proyecto aborda un problema de clasificación binaria: predecir si una afirmación política es verdadera o falsa a partir de su contenido textual e información contextual como el emisor, su afiliación política o el tema tratado.

Para resolver el problema se evaluaron distintos enfoques de aprendizaje automático, desde modelos clásicos hasta modelos más avanzados basados en \emph{embeddings} semánticos, técnicas de \emph{ensemble} y \emph{boosting}.

Esta memoria resume el análisis del conjunto de datos, el preprocesamiento realizado, los modelos desarrollados, la metodología de evaluación y las principales conclusiones obtenidas.

% ─────────────────────────────────────────────────────────────────────────────
\section{Descripción del problema}
% ─────────────────────────────────────────────────────────────────────────────

El objetivo del proyecto consiste en desarrollar un sistema capaz de clasificar afirmaciones políticas como verdaderas o falsas utilizando técnicas de Procesamiento del Lenguaje Natural y aprendizaje automático. Se trata de un problema de clasificación binaria, donde la variable objetivo toma el valor~1 para las afirmaciones verdaderas y~0 para las afirmaciones falsas.

El \emph{input} principal es el texto de la afirmación, pero también se cuenta con metadatos como tema, \emph{speaker}, ocupación, estado y afiliación política, que pueden aportar señal adicional al modelo.

El reto central es que la veracidad no siempre se puede inferir solo del texto. Dos afirmaciones pueden ser lingüísticamente similares y pertenecer a clases opuestas, por lo que no basta con un único enfoque. A eso se suma el \emph{class imbalance} del \emph{dataset}, razón por la que se usa \textbf{Macro F1} como métrica principal, ya que permite evaluar el modelo de forma balanceada entre ambas clases. Se probaron varios enfoques: desde modelos clásicos con representaciones tipo TF-IDF~\cite{sklearn} hasta modelos más avanzados con \emph{semantic embeddings}~\cite{sbert} y técnicas de \emph{ensemble}~\cite{lgbm,catboost,xgboost}, buscando el mejor \emph{trade-off} entre precisión y generalización.

% ─────────────────────────────────────────────────────────────────────────────
\section{Descripción del conjunto de datos}
% ─────────────────────────────────────────────────────────────────────────────

El conjunto de datos utilizado~\cite{liar} contiene afirmaciones políticas acompañadas de información contextual relacionada con cada una de ellas. Cada registro incluye el texto de la afirmación, recogido en la variable \texttt{statement}, junto con distintos metadatos como el tema tratado (\texttt{subject}), el hablante (\texttt{speaker}), su ocupación (\texttt{speaker\_job}), el estado (\texttt{state\_info}) y la afiliación política (\texttt{party\_affiliation}). La variable objetivo es \texttt{label}, donde el valor~1 representa una afirmación verdadera y el valor~0 una afirmación falsa.

\begin{table}[H]
  \centering
  \caption{Variables del \emph{dataset}.}
  \label{tab:columnas}
  \begin{tabularx}{\textwidth}{llX}
    \toprule
    \textbf{Variable}           & \textbf{Tipo}       & \textbf{Descripción} \\
    \midrule
    \texttt{id}                 & Identificador       & Solo para trazabilidad; excluido del modelo. \\
    \texttt{label}              & Target binario      & 0 = verdadero, 1 = falso. \\
    \texttt{statement}          & Texto               & El texto de la afirmación política. \\
    \texttt{subject}            & Categórica          & Tema(s) de la afirmación; puede contener varios. \\
    \texttt{speaker}            & Categórica          & Persona que realiza la afirmación. \\
    \texttt{speaker\_job}       & Categórica          & Cargo del hablante (27,7\,\% de valores ausentes). \\
    \texttt{state\_info}        & Categórica          & Estado geográfico de EE.\,UU. \\
    \texttt{party\_affiliation} & Categórica          & Afiliación política del hablante. \\
    \bottomrule
  \end{tabularx}
\end{table}

Durante el análisis exploratorio inicial se observó que el \emph{dataset} contiene \textbf{8.950 muestras} y presenta un cierto desbalanceo entre clases. Aproximadamente el~65\,\% de las muestras pertenecen a la clase verdadera, mientras que el~35\,\% corresponden a afirmaciones falsas. Este desequilibrio se tuvo en cuenta en la elección de la métrica de evaluación, dando prioridad a Macro F1, ya que permite valorar de forma más equilibrada el rendimiento en ambas clases.

También se detectaron valores nulos principalmente en las variables \texttt{speaker\_job} y \texttt{state\_info}. En lugar de eliminar estos registros, los valores faltantes se sustituyeron por la categoría \texttt{unknown}, conservando así todas las muestras disponibles y permitiendo que los modelos pudieran tratar la ausencia de información como una categoría más.

Otro aspecto relevante del \emph{dataset} es la alta cardinalidad de algunas variables categóricas, especialmente \texttt{speaker} y \texttt{speaker\_job}, donde muchos valores aparecen muy pocas veces. Para reducir ruido y mejorar la generalización, las categorías poco frecuentes se agruparon bajo la etiqueta \texttt{other}.

Por último, la columna \texttt{subject} podía contener varios temas separados por comas para una misma afirmación. Por ello, se realizó una limpieza de esta variable y se generaron características auxiliares, como una versión normalizada del texto y el número de temas asociados a cada afirmación.

% ─────────────────────────────────────────────────────────────────────────────
\section{Descripción de tareas comunes}
% ─────────────────────────────────────────────────────────────────────────────

El desarrollo del proyecto se llevó a cabo de forma colaborativa, distribuyendo las tareas entre los distintos integrantes del grupo. Cada miembro trabajó sobre diferentes etapas del \emph{pipeline} o sobre distintos modelos, pero las decisiones principales relacionadas con el preprocesamiento, la evaluación y la comparación de resultados se realizaron de forma conjunta para mantener coherencia en el trabajo.

La planificación del trabajo se estructuró en varias fases: análisis exploratorio del \emph{dataset}, limpieza y preprocesamiento de datos, generación de variables, entrenamiento y ajuste de modelos, evaluación de resultados y documentación final. Esta organización permitió comparar los distintos experimentos bajo criterios similares y seleccionar las estrategias con mejor rendimiento.

Para garantizar la reproducibilidad y el control de versiones, se utilizó GitHub~\cite{github} como repositorio común del proyecto. El desarrollo y la experimentación se realizaron principalmente en Jupyter Notebooks utilizando Python~\cite{python} y librerías de aprendizaje automático como \texttt{pandas}~\cite{pandas}, \texttt{scikit-learn}~\cite{sklearn}, XGBoost~\cite{xgboost}, LightGBM~\cite{lgbm}, CatBoost~\cite{catboost} y Hugging Face Transformers~\cite{hf_transformers}. La evaluación de resultados y generación de \emph{submissions} se llevó a cabo mediante Kaggle~\cite{kaggle}.

Para la búsqueda de hiperparámetros se utilizaron técnicas como \texttt{RandomizedSearchCV} y \texttt{HalvingGridSearchCV}~\cite{sklearn}, permitiendo explorar distintas configuraciones de forma eficiente. En los modelos Transformer también se ajustaron parámetros relacionados con el entrenamiento profundo, como \emph{learning rate}, \emph{batch size}, \emph{dropout} y número de épocas.

Finalmente, se compararon distintos enfoques de modelado, desde modelos lineales clásicos basados en TF-IDF hasta modelos \emph{ensemble} y arquitecturas Transformer avanzadas, analizando tanto el rendimiento obtenido como la capacidad de generalización de cada uno de ellos.

\subsection*{Herramientas y tecnologías}

\begin{table}[H]
  \centering
  \caption{Herramientas utilizadas en el proyecto.}
  \label{tab:tools}
  \renewcommand{\arraystretch}{1.25}
  \begin{tabularx}{\textwidth}{llX}
    \toprule
    \textbf{Categoría}         & \textbf{Herramienta}                                  & \textbf{Uso} \\
    \midrule
    Lenguaje                   & Python 3.x~\cite{python}                              & Todo el \emph{pipeline}. \\
    ML clásico                 & scikit-learn~\cite{sklearn}                           & LR, RFC, \texttt{StratifiedKFold}, \texttt{RandomizedSearchCV}, escalado, codificación. \\
    Gradient Boosting          & LightGBM~\cite{lgbm}                                  & Clasificador \emph{gradient boosting} hoja a hoja. \\
    Gradient Boosting          & CatBoost~\cite{catboost}                              & Soporte nativo de categóricas y \emph{ordered boosting}. \\
    Gradient Boosting          & XGBoost~\cite{xgboost}                                & Modelo base en el \emph{stacking}. \\
    \emph{Embeddings}          & sentence-transformers~\cite{sbert}                    & \texttt{all-MiniLM-L6-v2} (384~dim) y \texttt{all-mpnet-base-v2} (768~dim). \\
    NLP                        & NLTK~\cite{nltk}                                      & Porter/Snowball \emph{stemmer}, WordNet \emph{lemmatizer}. \\
    NER                        & spaCy~\cite{spacy}                                    & Conteo de entidades nombradas (PERSON, ORG, GPE, DATE\ldots). \\
    Transformers               & HuggingFace Transformers~\cite{hf_transformers}       & \emph{Fine-tuning} de DeBERTa-v3-small/base~\cite{deberta}. \\
    \emph{Fine-tuning} efic.   & PEFT + bitsandbytes~\cite{peft,bitsandbytes}          & LoRA sobre modelos 7B (experimento 5). \\
    \emph{Tracking}            & Weights \& Biases~\cite{wandb}                        & Métricas, curvas, importancia de \emph{features}, artefactos. \\
    Serialización              & joblib~\cite{joblib}                                  & Guardado de modelos, vectorizadores y \emph{threshold}. \\
    Control de versiones       & GitHub~\cite{github}                                  & Repositorio común del equipo. \\
    Evaluación                 & Kaggle~\cite{kaggle}                                  & Generación y envío de \emph{submissions}. \\
    \bottomrule
  \end{tabularx}
\end{table}

% ─────────────────────────────────────────────────────────────────────────────
\section{Evaluación}
% ─────────────────────────────────────────────────────────────────────────────

\subsection{Modelos evaluados}

Con el objetivo de comparar distintos enfoques para la detección de afirmaciones falsas, se evaluaron modelos de diferente complejidad, desde algoritmos lineales clásicos hasta modelos \emph{ensemble} y arquitecturas Transformer. La comparación permitió analizar tanto el rendimiento predictivo como la capacidad de generalización de cada técnica sobre el \emph{dataset}.

Los modelos fueron entrenados utilizando la misma estrategia de validación y conjuntos de datos procesados, adaptando únicamente la representación de las variables según las necesidades de cada algoritmo. Las métricas principales empleadas para la comparación fueron \textbf{Macro F1} y \textbf{ROC-AUC}, ya que el \emph{dataset} presentaba desbalanceo entre clases.

\subsubsection{Modelos lineales}

Los primeros modelos evaluados fueron algoritmos lineales clásicos basados en representaciones TF-IDF del texto~\cite{sklearn}. Se utilizaron principalmente Logistic Regression y Linear SVM, ya que suelen ofrecer buenos resultados como \emph{baseline} en problemas de clasificación textual.

Para estos modelos se emplearon representaciones TF-IDF con $n$-gramas de palabras y caracteres, además de variables categóricas relacionadas con el contexto político y el hablante. También se probaron distintas configuraciones de regularización y balanceo de clases mediante \texttt{class\_weight=\textquotedbl balanced\textquotedbl}.

La Regresión Logística destacó por su simplicidad, velocidad de entrenamiento e interpretabilidad, obteniendo resultados competitivos en las primeras fases del proyecto. Por otro lado, Linear SVM permitió mejorar ligeramente la separación entre clases, especialmente al incorporar $n$-gramas de caracteres y calibración de probabilidades.

Aunque estos modelos ofrecieron \emph{baselines} sólidos, presentaban limitaciones para capturar relaciones semánticas complejas entre el contenido textual y la información contextual.

\subsubsection{Modelos \emph{ensemble}}

Posteriormente se evaluaron distintos modelos \emph{ensemble} y basados en árboles de decisión, incluyendo Random Forest~\cite{sklearn}, XGBoost~\cite{xgboost}, LightGBM~\cite{lgbm} y CatBoost~\cite{catboost}. Estos modelos permitieron capturar relaciones no lineales entre las variables y aprovechar mejor las interacciones entre texto y \emph{metadata} contextual.

En Random Forest se utilizaron \emph{embeddings} semánticos~\cite{sbert} y variables derivadas relacionadas con el historial del hablante y el contexto político. Este modelo mostró una mejora significativa respecto a los modelos lineales, especialmente gracias al uso de \emph{embeddings} y variables agregadas.

XGBoost~\cite{xgboost} y LightGBM~\cite{lgbm} se emplearon junto con selección de características y ajuste de hiperparámetros mediante búsqueda aleatoria. Estos modelos consiguieron mejorar el rendimiento en ROC-AUC, aunque también mostraron sensibilidad al sobreajuste en determinadas variables relacionadas con el historial del hablante.

CatBoost~\cite{catboost} obtuvo uno de los mejores rendimientos entre los modelos individuales, gracias a su capacidad para manejar variables categóricas de forma nativa y reducir el impacto del \emph{leakage} mediante \emph{ordered boosting}.

Finalmente, se implementó una estrategia de \emph{stacking} combinando varios modelos base mediante un meta-clasificador Logistic Regression~\cite{sklearn}. Este enfoque permitió aprovechar la complementariedad entre modelos y mejorar tanto Macro F1 como ROC-AUC respecto a los modelos individuales.

\subsubsection{Transformers}

La última familia de modelos evaluados estuvo basada en arquitecturas Transformer preentrenadas~\cite{hf_transformers}, especialmente DeBERTa~\cite{deberta}. Estos modelos permiten capturar relaciones contextuales y semánticas más complejas que los enfoques basados únicamente en TF-IDF o \emph{embeddings} estáticos.

Para los Transformers se utilizó texto enriquecido con \emph{metadata} contextual, incorporando información como el hablante, el tema y la afiliación política directamente en la entrada textual del modelo. Además, se ajustaron hiperparámetros como \emph{learning rate}, \emph{batch size}, \emph{dropout} y número de épocas para optimizar el rendimiento. Para los experimentos de \emph{fine-tuning} eficiente se emplearon las librerías PEFT y bitsandbytes~\cite{peft,bitsandbytes}.

Los modelos Transformer mostraron mejores resultados que los modelos clásicos cuando se combinaron con información contextual. DeBERTa-v3-base con \emph{late fusion} en el \emph{stacking} alcanzó Macro F1 = 0,647 y ROC-AUC = 0,711.

Para superar ese techo se probó LoRA (\emph{Low-Rank Adaptation})~\cite{peft} sobre \texttt{mistralai/Mistral-7B-v0.1}~\cite{mistral}, cuantizado en 4 bits con bitsandbytes (NF4)~\cite{bitsandbytes}. Con \emph{k-fold} de 5 y \emph{learning rate} = 5e-5, el modelo alcanzó Macro F1 = \textbf{0,670} y ROC-AUC = \textbf{0,741} en modo \emph{standalone}, convirtiéndose en el mejor resultado del proyecto. Un hallazgo clave fue que añadir los GBDT al \emph{stacking} en este punto \emph{empeoró} el resultado: el coeficiente del transformer en el meta-clasificador ascendió a 2,75 frente a 0,08--0,50 de los modelos de árboles, evidencia de que estos ya no aportaban señal útil. El transformer con \emph{threshold sweep} directo sobre las predicciones OOF superó al \emph{stacking} completo en más de 0,05 puntos de Macro F1.

\subsection{Resultados}

La evaluación de los modelos se realizó utilizando validación cruzada estratificada de cinco \emph{folds} y un conjunto \emph{holdout} reservado para la evaluación final~\cite{sklearn}. Debido al desbalanceo presente en el \emph{dataset}, las métricas principales utilizadas fueron Macro F1 y ROC-AUC, ya que permiten evaluar de forma más equilibrada el rendimiento sobre ambas clases.

A lo largo del proyecto se observó una mejora progresiva del rendimiento al incorporar representaciones textuales más avanzadas, variables contextuales y técnicas de combinación de modelos. Los modelos lineales basados en TF-IDF proporcionaron \emph{baselines} sólidos, mientras que los modelos \emph{ensemble} y Transformer consiguieron capturar relaciones más complejas entre el texto y la \emph{metadata} política.

Los resultados finales obtenidos por los modelos más relevantes se muestran en la Tabla~\ref{tab:resultados}.

\begin{table}[H]
  \centering
  \caption{Resultados en \emph{holdout} de los modelos evaluados (ordenados por Macro F1).}
  \label{tab:resultados}
  \renewcommand{\arraystretch}{1.3}
  \begin{tabular}{lccc}
    \toprule
    \textbf{Modelo}                                                        & \textbf{Macro F1} & \textbf{ROC-AUC} & \textbf{Umbral} \\
    \midrule
    LR \emph{baseline} (TF-IDF)                                            & 0,604             & 0,654            & 0,50 \\
    LR + \emph{nested CV} + \emph{embeddings}                              & 0,611             & 0,660            & ---  \\
    Random Forest (embeddings + NER + \texttt{job\_true\_rate})            & 0,621             & 0,667            & 0,58 \\
    LGBM Opción B (sin \texttt{speaker\_true\_rate})                        & 0,618             & 0,679            & 0,52 \\
    CatBoost + Opción B                                                     & 0,629             & 0,674            & 0,46 \\
    \emph{Stacking} (\texttt{all-mpnet-base-v2}, 768\,dim)                  & 0,643             & 0,700            & 0,62 \\
    \emph{Late Fusion} + DeBERTa-v3-small                                   & 0,644             & 0,709            & ---  \\
    \emph{Late Fusion} + DeBERTa-v3-base                                    & 0,647             & 0,711            & ---  \\
    Mistral-7B LoRA \emph{standalone} (fold\,1 colapsado)                   & 0,655             & 0,724            & 0,50 \\
    Mistral-7B LoRA + \emph{Late Fusion} (Lora\,1)                          & 0,649             & 0,721            & 0,63 \\
    Mistral-7B LoRA estable + \emph{Stacking} (Lora\,2)                     & 0,664             & 0,742            & 0,57 \\
    \textbf{Mistral-7B LoRA estable \emph{standalone} (Lora\,2)}            & \textbf{0,670}    & \textbf{0,741}   & \textbf{0,41} \\
    Mistral-7B LoRA 3\,épocas \emph{standalone} (Lora\,3)                   & 0,668             & 0,748            & 0,43 \\
    Mistral-7B LoRA 3\,épocas \emph{transformer-only} (Lora\,3)             & 0,669             & 0,741            & 0,43 \\
    \bottomrule
  \end{tabular}
\end{table}

Uno de los factores que más contribuyó a la mejora del rendimiento fue la incorporación de \emph{embeddings} semánticos~\cite{sbert} y variables contextuales relacionadas con el hablante y el tema de la afirmación. También se comprobó que algunas variables históricas, aunque muy predictivas, podían introducir problemas de sobreajuste si no se calculaban correctamente mediante estrategias \emph{out-of-fold}.

Los modelos Transformer~\cite{hf_transformers,deberta} mejoraron notablemente al incorporar \emph{metadata} contextual directamente en el texto de entrada. Hasta DeBERTa-v3-base, los mejores resultados se obtenían mediante \emph{late fusion} combinando Transformers con modelos \emph{ensemble}. Sin embargo, al escalar a Mistral-7B con LoRA~\cite{peft,mistral}, el patrón se invirtió: el modelo de lenguaje capturaba señal de texto tan diferente a la de los GBDT que añadir estos al \emph{ensemble} introducía ruido en lugar de complementariedad. El transformer \emph{standalone} superó al \emph{stacking} completo en Macro F1.

En general, los resultados muestran que la combinación de representaciones y enfoques es beneficiosa hasta el punto en que la diferencia de calidad entre los componentes es demasiado grande; en ese momento, el componente más fuerte debe utilizarse solo.

% ─────────────────────────────────────────────────────────────────────────────
\section{Análisis y conclusiones conjuntas}
% ─────────────────────────────────────────────────────────────────────────────

\subsection{Análisis del preprocesamiento y la generación de variables}

El preprocesamiento de datos y la generación de variables tuvieron un papel fundamental en el rendimiento de los modelos desarrollados. El objetivo principal fue reducir el ruido del \emph{dataset}, representar correctamente la información textual y contextual, y evitar problemas de sobreajuste o \emph{leakage} durante el entrenamiento.

En primer lugar, se realizó una limpieza general de las variables categóricas y textuales. Las columnas de texto se normalizaron convirtiendo todo el contenido a minúsculas, eliminando espacios innecesarios y estandarizando formatos inconsistentes. Los valores nulos de variables como \texttt{speaker\_job} y \texttt{state\_info} se reemplazaron por categorías específicas como \texttt{unknown} o \texttt{missing}, evitando la pérdida de información asociada a la ausencia de datos.

Para los modelos clásicos basados en TF-IDF~\cite{sklearn} se aplicó un preprocesamiento textual más agresivo sobre la columna \texttt{statement}. Este proceso incluyó eliminación de puntuación, normalización de contracciones, reemplazo de números por un token genérico, tokenización, eliminación de \emph{stopwords} y lematización con NLTK~\cite{nltk}. Además, se conservaron negaciones como \emph{not} o \emph{never}, ya que aportaban información relevante para la clasificación.

La representación textual se realizó principalmente mediante TF-IDF con $n$-gramas de palabras y caracteres~\cite{sklearn}. También se utilizaron \emph{embeddings} semánticos generados con Sentence Transformers~\cite{sbert}, como \texttt{all-MiniLM-L6-v2} y \texttt{all-mpnet-base-v2}, permitiendo capturar relaciones semánticas más complejas entre afirmaciones.

En las variables categóricas se aplicaron técnicas de reducción de cardinalidad agrupando categorías poco frecuentes bajo la etiqueta \texttt{other}. Esto se utilizó especialmente en variables como \texttt{speaker}, \texttt{speaker\_job} y \texttt{subject}, donde existía una gran cantidad de valores con muy pocas apariciones.

La variable \texttt{subject} requería un tratamiento especial debido a que una misma afirmación podía pertenecer a varios temas simultáneamente. Para ello se utilizó \texttt{MultiLabelBinarizer}~\cite{sklearn}, generando variables binarias independientes para cada tema relevante.

Además, se generaron variables derivadas relacionadas con el contexto político y el comportamiento histórico de los hablantes. Entre ellas destacan las tasas históricas de veracidad asociadas al hablante, partido político o tema tratado. Para evitar \emph{leakage}, estas estadísticas se calcularon utilizando estrategias \emph{out-of-fold} dentro de la validación cruzada~\cite{sklearn}.

Finalmente, en los modelos Transformer~\cite{hf_transformers} no se aplicó una limpieza textual agresiva, ya que estos modelos aprovechan mejor el lenguaje natural original. En su lugar, se construyó un texto enriquecido que combinaba la afirmación con información contextual como el hablante, el tema y la afiliación política, permitiendo incorporar \emph{metadata} directamente en la entrada del modelo.

\subsection{Conclusiones}

En este proyecto se abordó un problema de clasificación de afirmaciones políticas~\cite{liar} mediante técnicas de procesamiento del lenguaje natural y aprendizaje automático. A lo largo del desarrollo se evaluaron distintos enfoques, desde modelos lineales clásicos hasta arquitecturas Transformer~\cite{hf_transformers,deberta} y estrategias \emph{ensemble} avanzadas~\cite{lgbm,catboost,xgboost}.

Los resultados obtenidos muestran que la información contextual asociada a las afirmaciones, como el hablante, el tema o la afiliación política, tuvo un impacto muy relevante en el rendimiento de los modelos. La combinación de esta \emph{metadata} con representaciones textuales avanzadas~\cite{sbert} permitió mejorar significativamente las métricas de clasificación.

Los modelos lineales basados en TF-IDF~\cite{sklearn} ofrecieron \emph{baselines} sólidos y eficientes, mientras que los modelos \emph{ensemble}~\cite{lgbm,catboost,xgboost} lograron capturar relaciones no lineales entre variables y mejorar la capacidad predictiva. Por su parte, los modelos Transformer~\cite{hf_transformers,deberta} consiguieron representar mejor el contexto semántico del lenguaje natural, especialmente al incorporar información contextual directamente en el texto de entrada.

También se observó la importancia de controlar el sobreajuste y evitar \emph{leakage} en variables relacionadas con el historial de veracidad de los hablantes. El uso de validación cruzada estratificada~\cite{sklearn} y estrategias \emph{out-of-fold} permitió reducir este problema y obtener estimaciones más realistas del rendimiento.

El mejor resultado del proyecto se obtuvo con Mistral-7B LoRA~\cite{peft,mistral} en modo \emph{standalone}, alcanzando Macro F1 = 0,670 y ROC-AUC = 0,741. Una lección clave de la fase final fue que el \emph{stacking} con GBDT resultó contraproducente cuando el transformer dominaba el \emph{ensemble}: el coeficiente del meta-clasificador ascendió a 2,75 para el transformer frente a 0,08--0,50 para los modelos de árboles. Eliminar los GBDT y aplicar directamente un \emph{threshold sweep} sobre las predicciones OOF del transformer mejoró el Macro F1 en más de 0,05 puntos sobre la \emph{baseline} de \emph{stacking}. El \emph{ensemble} es una herramienta útil cuando los modelos son complementarios, pero resulta perjudicial cuando uno de ellos es cualitativamente superior.

También se comprobó la importancia de la tasa de aprendizaje en el \emph{fine-tuning} con LoRA: con LR\,=\,2e-4, el \emph{score head} inicializado aleatoriamente generaba pérdidas extremas en el \emph{batch}\,0 que impedían la convergencia de un \emph{fold} completo. Reducir a LR\,=\,5e-5 e inicializar el \emph{score head} con desviación estándar 0,01 dio cinco \emph{folds} estables y el mejor Macro F1 del proyecto.

En conclusión, el proyecto permitió comparar técnicas de clasificación textual desde modelos lineales hasta LLM de 7B parámetros, comprobando cómo la integración de información contextual, el ajuste cuidadoso de hiperparámetros y la elección correcta de la estrategia de \emph{ensemble} son determinantes para el rendimiento en tareas de detección de desinformación política.

% ─────────────────────────────────────────────────────────────────────────────
\bibliographystyle{unsrtnat}
\bibliography{references}
% ─────────────────────────────────────────────────────────────────────────────

\end{document}
